% ============================================================
% SECTIONS MANQUANTES À INTÉGRER DANS LE RAPPORT PFE
% Farm AI — Éléments identifiés après analyse du code source
% ============================================================

% ============================================================
% 1. REMPLACER la section "Présentation des acteurs" (Chapitre 2)
%    par cette version corrigée avec les 4 rôles
% ============================================================

\subsection{Présentation des acteurs}

Le système Farm AI repose sur un modèle multi-rôles avec quatre types d'acteurs distincts,
chacun disposant d'un niveau d'accès et de responsabilités spécifiques :

\begin{itemize}
  \item \textbf{Administrateur (ROLE\_ADMIN) :} acteur principal du système. Il dispose d'un
    accès complet à l'ensemble des fonctionnalités : gestion des employés, des gestionnaires,
    des observateurs, des départements, des caméras, consultation du tableau de bord analytique,
    gestion des alertes et configuration de son compte. Il peut s'authentifier par email/mot de
    passe ou par reconnaissance faciale biométrique.
  \item \textbf{Gestionnaire (ROLE\_MANAGER) :} responsable d'un département spécifique. Il
    accède à un tableau de bord dédié affichant les employés assignés à son département, les
    caméras associées et les alertes. Il peut valider ou modifier la composition de son équipe
    via un système de suggestion automatique. Son compte est créé par l'administrateur et
    activé par email.
  \item \textbf{Observateur (ROLE\_VIEWER) :} acteur disposant d'un accès en lecture seule.
    Il consulte la liste des employés enregistrés, les départements configurés et peut
    déclencher l'enregistrement facial d'un employé. Son compte est créé par l'administrateur
    et activé par email.
  \item \textbf{Service IA Python :} acteur secondaire automatisé. Une fois activé, il traite
    les flux vidéo de manière autonome (détection YOLO, reconnaissance faciale FaceNet) et
    envoie les résultats au backend sans intervention humaine.
\end{itemize}

% ============================================================
% 2. TABLEAU DES BESOINS FONCTIONNELS ET NON FONCTIONNELS
%    Remplacer les listes à puces existantes
% ============================================================

\subsection{Besoins fonctionnels et non fonctionnels}

Le tableau~\ref{tab:besoins_fonctionnels} récapitule l'ensemble des besoins fonctionnels
identifiés, classés par module et par acteur concerné.

\begin{longtable}{|c|>{\raggedright\arraybackslash}p{5cm}|>{\raggedright\arraybackslash}p{6cm}|l|}
  \caption{Besoins fonctionnels du système Farm AI} \label{tab:besoins_fonctionnels} \\
  \hline
  \rowcolor{gray!20}
  \textbf{ID} & \textbf{Besoin fonctionnel} & \textbf{Description} & \textbf{Acteur} \\
  \hline
  \endfirsthead
  \hline
  \rowcolor{gray!20}
  \textbf{ID} & \textbf{Besoin fonctionnel} & \textbf{Description} & \textbf{Acteur} \\
  \hline
  \endhead
  \hline
  \endfoot

  BF01 & Authentification classique &
    Se connecter via email et mot de passe sécurisé avec génération de token JWT. &
    Tous \\
  \hline
  BF02 & Authentification biométrique &
    Se connecter par reconnaissance faciale en temps réel sans saisie de mot de passe. &
    Tous \\
  \hline
  BF03 & Gérer les employés &
    Ajouter, modifier, supprimer et consulter les fiches des employés avec leur métier, email et statut. &
    Admin \\
  \hline
  BF04 & Enregistrer un Face ID employé &
    Capturer et enregistrer l'encodage biométrique du visage d'un employé pour le pointage automatique. &
    Admin \\
  \hline
  BF05 & Approuver/Rejeter un employé &
    Valider ou refuser un employé nouvellement créé avant son affectation à un département. &
    Admin \\
  \hline
  BF06 & Gérer les gestionnaires &
    Créer, modifier et supprimer des comptes gestionnaires. Un email d'activation avec mot de passe est envoyé automatiquement via Gmail. &
    Admin \\
  \hline
  BF07 & Gérer les observateurs &
    Créer, modifier et supprimer des comptes observateurs. Un email d'activation est envoyé automatiquement via Gmail. &
    Admin \\
  \hline
  BF08 & Gérer les départements &
    Créer, modifier et supprimer des départements avec nom, horaires, responsable (manager) et besoins en personnel par métier. &
    Admin \\
  \hline
  BF09 & Activer un compte par email &
    Cliquer sur le lien d'activation reçu par Gmail pour activer son compte et enregistrer son visage. &
    Manager / Viewer \\
  \hline
  BF10 & Enregistrer son visage (onboarding) &
    Enregistrer ses données biométriques lors de la première connexion après activation du compte. &
    Manager / Viewer \\
  \hline
  BF11 & Gérer les caméras &
    Consulter la liste des caméras, vérifier leur statut (active/inactive) et déclencher l'analyse IA. &
    Admin \\
  \hline
  BF12 & Lancer la surveillance IA &
    Activer l'analyse intelligente sur une caméra : détection de personnes (YOLO), reconnaissance faciale (FaceNet) et analyse vestimentaire. &
    Admin \\
  \hline
  BF13 & Pointage automatique &
    Enregistrer automatiquement les entrées et sorties des employés identifiés avec horodatage précis. &
    Service IA \\
  \hline
  BF14 & Consulter le tableau de bord &
    Visualiser les KPIs (présents, heures travaillées, entrées, caméras actives), les graphiques de présence et les insights IA. &
    Admin \\
  \hline
  BF15 & Gérer les alertes &
    Consulter les alertes générées lors de détections d'anomalies (personne non identifiée, intrusion). &
    Admin \\
  \hline
  BF16 & Suggestion d'équipe &
    Recevoir une proposition automatique d'équipe basée sur les besoins en personnel du département et la disponibilité des employés. &
    Manager \\
  \hline
  BF17 & Valider/Modifier l'équipe &
    Accepter la suggestion d'équipe ou la modifier manuellement avant validation. Un email de confirmation est envoyé à l'admin et aux employés. &
    Manager \\
  \hline
  BF18 & Tableau de bord manager &
    Consulter un tableau de bord dédié avec les employés du département, les caméras associées et les alertes. &
    Manager \\
  \hline
  BF19 & Tableau de bord observateur &
    Consulter en lecture seule la liste des employés et des départements. &
    Viewer \\
  \hline
  BF20 & Notifications par email (Gmail) &
    Envoyer automatiquement des emails via Gmail SMTP pour : création de compte, activation, affectation, validation d'équipe et suppression. &
    Système \\
  \hline
  BF21 & Gérer son compte &
    Modifier ses informations personnelles, changer son mot de passe et gérer ses données biométriques. &
    Tous \\
  \hline
  BF22 & Changer la langue &
    Basculer l'interface entre le français et l'anglais. &
    Tous \\
  \hline
\end{longtable}

Le tableau~\ref{tab:besoins_non_fonctionnels} présente les besoins non fonctionnels qui ont
guidé les choix techniques tout au long du développement.

\begin{table}[H]
  \centering
  \caption{Besoins non fonctionnels du système Farm AI}
  \label{tab:besoins_non_fonctionnels}
  \renewcommand{\arraystretch}{1.4}
  \begin{tabularx}{\textwidth}{|c|l|X|}
    \hline
    \rowcolor{gray!20}
    \textbf{ID} & \textbf{Catégorie} & \textbf{Description} \\
    \hline
    BNF01 & Performance &
      Le traitement du flux vidéo doit s'effectuer avec une latence inférieure à 2 secondes entre la capture et l'affichage du résultat. \\
    \hline
    BNF02 & Sécurité &
      Toutes les requêtes sont authentifiées par JWT. Les données biométriques sont stockées sous forme d'encodages mathématiques. Les mots de passe sont hachés avec BCrypt. \\
    \hline
    BNF03 & Fiabilité IA &
      Le taux d'identification correcte doit être supérieur à 90\,\%. Le modèle FaceNet retenu atteint 97,2\,\%. \\
    \hline
    BNF04 & Ergonomie &
      L'interface doit être intuitive pour un utilisateur non technicien. Toutes les actions critiques sont confirmées par une fenêtre de validation. \\
    \hline
    BNF05 & Évolutivité &
      L'architecture micro-services permet de faire évoluer chaque composant indépendamment. \\
    \hline
    BNF06 & Contrôle d'accès &
      Le système applique un contrôle d'accès basé sur les rôles (RBAC) avec redirection automatique selon le rôle de l'utilisateur connecté. \\
    \hline
    BNF07 & Disponibilité &
      Le service de notification par email doit fonctionner de manière asynchrone pour ne pas bloquer les opérations principales. \\
    \hline
    BNF08 & Internationalisation &
      L'interface supporte le français et l'anglais avec basculement dynamique sans rechargement. \\
    \hline
    BNF09 & Compatibilité &
      L'application est compatible avec les navigateurs modernes (Chrome, Firefox, Edge). \\
    \hline
  \end{tabularx}
\end{table}

% ============================================================
% 3. SECTION GMAIL / NOTIFICATIONS PAR EMAIL
%    À ajouter dans le Chapitre 2, section Environnement logiciel
% ============================================================

\subsection{Service de notification par email (Gmail SMTP)}

Le système Farm AI intègre un service de notification par email utilisant le protocole
\textbf{Gmail SMTP}. Ce service, implémenté dans le backend Spring Boot via la classe
\texttt{EmailService}, utilise l'API \texttt{JavaMailSender} de Spring pour envoyer des
emails de manière asynchrone (annotation \texttt{@Async}).

La configuration dans le fichier \texttt{application.properties} utilise le serveur
SMTP de Google (\texttt{smtp.gmail.com}, port 587) avec l'authentification TLS activée.
Un mot de passe d'application Google est utilisé pour sécuriser la connexion.

Le service gère huit types de notifications automatiques :

\begin{enumerate}
  \item \textbf{Création d'un compte employé :} notification informant l'employé que son
    profil a été créé dans le système.
  \item \textbf{Validation d'un employé :} email de bienvenue envoyé après approbation par
    l'administrateur.
  \item \textbf{Activation de compte gestionnaire :} email contenant le mot de passe
    généré et un lien d'activation unique avec token.
  \item \textbf{Activation de compte observateur :} même mécanisme que pour les gestionnaires,
    avec un lien d'activation dédié.
  \item \textbf{Affectation à un département :} notification envoyée à l'employé lors de
    son assignation, incluant le nom du département, le responsable et les horaires.
  \item \textbf{Validation d'équipe (vers l'admin) :} récapitulatif envoyé à l'administrateur
    lorsqu'un gestionnaire valide la composition de son équipe.
  \item \textbf{Confirmation d'équipe (vers le manager) :} confirmation envoyée au
    gestionnaire après validation réussie.
  \item \textbf{Suppression de département :} notification informant le gestionnaire
    de la suppression de son département.
\end{enumerate}

% ============================================================
% 4. SECTION GESTION DES GESTIONNAIRES (Managers)
%    À ajouter comme spécification fonctionnelle
% ============================================================

\subsection{Gérer les gestionnaires}

\begin{table}[H]
  \centering
  \caption{Description textuelle du cas d'utilisation « Gérer les gestionnaires »}
  \label{tab:cu_managers}
  \renewcommand{\arraystretch}{1.4}
  \begin{tabularx}{\textwidth}{|>{\bfseries}l|L|}
    \hline
    \rowcolor{gray!20}
    \multicolumn{2}{|c|}{\textbf{Créer un gestionnaire}} \\
    \hline
    Nom C.U & Créer un gestionnaire \\
    \hline
    Acteur & Administrateur \\
    \hline
    Description & L'administrateur crée un compte gestionnaire. Un mot de passe est généré
      automatiquement et un email d'activation est envoyé via Gmail. \\
    \hline
    Précondition & L'administrateur est authentifié et se trouve sur le module de gestion
      des gestionnaires. \\
    \hline
    Postcondition & Le compte gestionnaire est créé en base de données avec le rôle
      ROLE\_MANAGER, le statut désactivé, et un email d'activation est envoyé. \\
    \hline
    Scénario nominal &
      1. L'administrateur saisit le prénom, nom, email et téléphone.\newline
      2. Il clique sur « Enregistrer ».\newline
      3. Le backend génère un mot de passe aléatoire et un token d'activation.\newline
      4. Le compte est créé avec \texttt{enabled = false}.\newline
      5. Un email contenant le mot de passe et le lien d'activation est envoyé via Gmail SMTP.\newline
      6. Le gestionnaire apparaît dans la liste avec le statut « En attente d'activation ». \\
    \hline
    Scénarios alternatifs &
      -- Si l'email existe déjà, une erreur 400 est retournée.\newline
      -- Si le service email est indisponible, le compte est créé mais l'email n'est pas envoyé. \\
    \hline
  \end{tabularx}
\end{table}

% ============================================================
% 5. SECTION GESTION DES DÉPARTEMENTS
% ============================================================

\subsection{Gérer les départements}

\begin{table}[H]
  \centering
  \caption{Description textuelle du cas d'utilisation « Créer un département »}
  \label{tab:cu_departments}
  \renewcommand{\arraystretch}{1.4}
  \begin{tabularx}{\textwidth}{|>{\bfseries}l|L|}
    \hline
    \rowcolor{gray!20}
    \multicolumn{2}{|c|}{\textbf{Créer un département}} \\
    \hline
    Nom C.U & Créer un département \\
    \hline
    Acteur & Administrateur \\
    \hline
    Description & L'administrateur crée un département en définissant son nom, ses horaires
      de travail, le gestionnaire responsable et les besoins en personnel par métier. \\
    \hline
    Précondition & L'administrateur est authentifié. Au moins un gestionnaire existe
      dans le système. \\
    \hline
    Postcondition & Le département est enregistré en base avec les besoins en docteurs,
      électriciens et ouvriers. \\
    \hline
    Scénario nominal &
      1. L'administrateur saisit le nom du département.\newline
      2. Il sélectionne un gestionnaire responsable dans la liste déroulante.\newline
      3. Il définit les horaires d'ouverture et de fermeture.\newline
      4. Il précise les besoins en personnel : nombre de docteurs, d'électriciens et d'ouvriers requis.\newline
      5. Il clique sur « Enregistrer ».\newline
      6. Le département est créé et affiché sous forme de carte avec les détails. \\
    \hline
    Scénarios alternatifs &
      -- Si aucun gestionnaire n'est sélectionné, un avertissement s'affiche.\newline
      -- Si les nombres d'employés sont négatifs, une validation empêche la soumission. \\
    \hline
  \end{tabularx}
\end{table}

% ============================================================
% 6. SECTION ACTIVATION DE COMPTE ET FACE SETUP
% ============================================================

\subsection{Activer un compte et enregistrer son visage}

\begin{table}[H]
  \centering
  \caption{Description textuelle du cas d'utilisation « Activer un compte »}
  \label{tab:cu_activate}
  \renewcommand{\arraystretch}{1.4}
  \begin{tabularx}{\textwidth}{|>{\bfseries}l|L|}
    \hline
    \rowcolor{gray!20}
    \multicolumn{2}{|c|}{\textbf{Activer un compte par email}} \\
    \hline
    Nom C.U & Activer un compte par email \\
    \hline
    Acteur & Gestionnaire / Observateur \\
    \hline
    Description & Le gestionnaire ou l'observateur clique sur le lien d'activation reçu
      par email pour activer son compte, puis enregistre son visage. \\
    \hline
    Précondition & Un email d'activation a été envoyé. Le token est valide et non expiré. \\
    \hline
    Postcondition & Le compte est activé (\texttt{enabled = true}) et l'utilisateur est
      redirigé vers la page d'enregistrement facial. \\
    \hline
    Scénario nominal &
      1. L'utilisateur clique sur le lien d'activation dans l'email reçu.\newline
      2. L'application Angular extrait le token de l'URL.\newline
      3. Une requête POST est envoyée au backend pour valider le token.\newline
      4. Le backend active le compte et retourne un token JWT.\newline
      5. L'utilisateur est redirigé vers la page « Face Setup ».\newline
      6. Il enregistre son visage via la caméra.\newline
      7. Après succès, il est redirigé vers son tableau de bord dédié. \\
    \hline
    Scénarios alternatifs &
      -- Si le token est invalide ou expiré, un message d'erreur s'affiche.\newline
      -- Si la caméra est indisponible, l'enregistrement facial échoue. \\
    \hline
  \end{tabularx}
\end{table}

% ============================================================
% 7. SECTION SUGGESTION ET VALIDATION D'ÉQUIPE
% ============================================================

\subsection{Suggestion et validation d'équipe}

\begin{table}[H]
  \centering
  \caption{Description textuelle du cas d'utilisation « Valider une équipe »}
  \label{tab:cu_team}
  \renewcommand{\arraystretch}{1.4}
  \begin{tabularx}{\textwidth}{|>{\bfseries}l|L|}
    \hline
    \rowcolor{gray!20}
    \multicolumn{2}{|c|}{\textbf{Valider une équipe}} \\
    \hline
    Nom C.U & Valider une équipe \\
    \hline
    Acteur & Gestionnaire \\
    \hline
    Description & Le gestionnaire consulte la proposition automatique d'équipe générée
      par le système, la modifie si nécessaire, puis la valide pour affecter les employés
      à son département. \\
    \hline
    Précondition & Le gestionnaire est authentifié. Son département a des besoins en
      personnel définis. Des employés disponibles et approuvés existent en base. \\
    \hline
    Postcondition & Les employés sélectionnés sont affectés au département. Des emails
      de confirmation sont envoyés au gestionnaire, à l'administrateur et aux employés. \\
    \hline
    Scénario nominal &
      1. Le gestionnaire accède à son tableau de bord.\newline
      2. Le système affiche une suggestion d'équipe basée sur les besoins du département
         et la disponibilité des employés (visage enregistré, statut APPROVED, non assigné).\newline
      3. Le gestionnaire peut ajouter ou retirer des employés de la sélection.\newline
      4. Il clique sur « Valider l'équipe ».\newline
      5. Le backend vérifie que chaque employé est disponible et validé.\newline
      6. Les employés sont assignés au département (\texttt{available = false}).\newline
      7. Un email de confirmation est envoyé au gestionnaire, à l'admin et à chaque employé. \\
    \hline
    Scénarios alternatifs &
      -- Si un employé n'est pas approuvé ou n'a pas de visage enregistré, la validation échoue.\newline
      -- Le gestionnaire peut modifier l'équipe après validation initiale. \\
    \hline
  \end{tabularx}
\end{table}

% ============================================================
% 8. ENTRÉES DU BACKLOG MANQUANTES
%    À ajouter au tableau du Product Backlog
% ============================================================

% Lignes à AJOUTER dans le tableau \ref{tab:backlog} :
%
%    Gérer les gestionnaires &
%      En tant qu'administrateur, je peux créer, modifier et supprimer des comptes
%      gestionnaires. Un email d'activation est envoyé automatiquement via Gmail. &
%      Administrateur \\
%    \hline
%    Gérer les observateurs &
%      En tant qu'administrateur, je peux créer, modifier et supprimer des comptes
%      observateurs en lecture seule. Un email d'activation est envoyé via Gmail. &
%      Administrateur \\
%    \hline
%    Gérer les départements &
%      En tant qu'administrateur, je peux créer des départements avec horaires,
%      responsable et besoins en personnel par métier. &
%      Administrateur \\
%    \hline
%    Activer un compte &
%      En tant que gestionnaire/observateur, je peux activer mon compte en cliquant
%      sur le lien reçu par email et en enregistrant mon visage. &
%      Manager / Viewer \\
%    \hline
%    Suggestion d'équipe &
%      En tant que gestionnaire, je reçois une proposition automatique d'équipe et
%      je peux la valider ou la modifier. &
%      Gestionnaire \\
%    \hline
%    Tableau de bord manager &
%      En tant que gestionnaire, je peux consulter un tableau de bord dédié à mon département. &
%      Gestionnaire \\
%    \hline
%    Tableau de bord observateur &
%      En tant qu'observateur, je peux consulter en lecture seule la liste des employés
%      et des départements. &
%      Observateur \\
%    \hline
%    Notifications Gmail &
%      En tant que système, j'envoie automatiquement des emails via Gmail SMTP pour
%      chaque événement important. &
%      Système \\
%    \hline

% ============================================================
% 9. CORRECTION : La base de données est MySQL, pas PostgreSQL
%    Remplacer toute mention de PostgreSQL par MySQL
% ============================================================

% AVANT : PostgreSQL est la base de données relationnelle utilisée...
% APRÈS :
\subsection{MySQL (correction)}

MySQL est la base de données relationnelle utilisée pour la persistance de toutes
les données du système : utilisateurs (administrateurs, gestionnaires, observateurs),
employés, départements, pointages, caméras et alertes. La connexion est configurée
dans le fichier \texttt{application.properties} avec le driver JDBC MySQL sur le
port 3306, et Hibernate gère automatiquement le schéma via \texttt{ddl-auto=update}.

% ============================================================
% 10. INTERFACES MANQUANTES À DOCUMENTER
%     Ajouter dans les sections Réalisation des sprints concernés
% ============================================================

% --- Interface de gestion des gestionnaires ---
% \begin{figure}[H]
%   \centering
%   \includegraphics[width=0.90\textwidth]{images/managers.png}
%   \caption{Interface de gestion des gestionnaires}
%   \label{fig:managers}
% \end{figure}
%
% Description : Cette interface permet à l'administrateur de créer, modifier et supprimer
% des comptes gestionnaires. Le formulaire collecte le prénom, nom, email et téléphone.
% La liste affiche le statut d'activation (ACTIF / EN ATTENTE) et l'état de
% l'enregistrement facial (VISAGE OK / VISAGE MANQUANT) pour chaque gestionnaire.

% --- Interface de gestion des observateurs ---
% \begin{figure}[H]
%   \centering
%   \includegraphics[width=0.90\textwidth]{images/viewers.png}
%   \caption{Interface de gestion des observateurs}
%   \label{fig:viewers}
% \end{figure}

% --- Interface de gestion des départements ---
% \begin{figure}[H]
%   \centering
%   \includegraphics[width=0.90\textwidth]{images/departments.png}
%   \caption{Interface de gestion des départements}
%   \label{fig:departments}
% \end{figure}
%
% Description : Cette interface permet de créer des départements en définissant le nom,
% les horaires de travail (heure d'ouverture/fermeture), le gestionnaire responsable et
% les besoins exacts en personnel par métier (docteurs, électriciens, ouvriers). Les
% départements sont affichés sous forme de cartes avec le nombre d'employés affectés.

% --- Interface du tableau de bord gestionnaire ---
% \begin{figure}[H]
%   \centering
%   \includegraphics[width=0.90\textwidth]{images/manager_dashboard.png}
%   \caption{Tableau de bord du gestionnaire}
%   \label{fig:manager_dashboard}
% \end{figure}
%
% Description : Ce tableau de bord dédié affiche les informations du département du
% gestionnaire : employés assignés, caméras associées, alertes actives. Le système de
% suggestion d'équipe propose automatiquement des employés disponibles, que le gestionnaire
% peut accepter, modifier ou valider.

% --- Interface d'activation de compte ---
% \begin{figure}[H]
%   \centering
%   \includegraphics[width=0.90\textwidth]{images/activate.png}
%   \caption{Page d'activation de compte}
%   \label{fig:activate}
% \end{figure}

% --- Interface d'enregistrement facial (Face Setup) ---
% \begin{figure}[H]
%   \centering
%   \includegraphics[width=0.90\textwidth]{images/face_setup.png}
%   \caption{Page d'enregistrement facial (onboarding)}
%   \label{fig:face_setup}
% \end{figure}

% ============================================================
% 11. PLANIFICATION DES SPRINTS MISE À JOUR
%     Remplacer le tableau existant
% ============================================================

\begin{table}[H]
  \centering
  \caption{Planification des sprints (mise à jour)}
  \label{tab:sprints_v2}
  \renewcommand{\arraystretch}{1.5}
  \begin{tabularx}{\textwidth}{|c|X|c|}
    \hline
    \rowcolor{gray!20}
    \textbf{Sprint} & \textbf{Contenu} & \textbf{Durée} \\
    \hline
    Sprint 0 & Spécification des besoins, backlog produit, planification, environnement technique, architecture & 1 semaine \\
    \hline
    Sprint 1 & Authentification (classique et biométrique), gestion des employés, enregistrement Face ID, gestion des gestionnaires et observateurs, activation de compte par email, notifications Gmail & 2 semaines \\
    \hline
    Sprint 2 & Gestion des départements, gestion des caméras, surveillance IA, pointage automatique, centre d'alertes, suggestion et validation d'équipe & 2 semaines \\
    \hline
    Sprint 3 & Tableau de bord analytique (admin), tableau de bord gestionnaire, tableau de bord observateur, paramètres du compte, tests de validation & 2 semaines \\
    \hline
  \end{tabularx}
\end{table}

\end{document}
